\chapter{eBPF em Containers e Kubernetes}
\label{cap:containers}

A adoção massiva de contêineres e orquestradores como o Kubernetes trouxe novos
desafios de segurança: ambientes altamente dinâmicos, com processos efêmeros,
múltiplos inquilinos compartilhando o mesmo kernel e superfície de ataque
distribuída por dezenas ou centenas de nós \cite{cassagnes2020}.
Abordagens tradicionais de segurança baseadas em agentes por processo ou em
inspeção de logs com atraso mostram-se inadequadas nesse contexto.
O eBPF, por operar no kernel compartilhado por todos os contêineres de um nó,
oferece visibilidade e controle centralizados sem depender de instrumentação
individual de cada carga de trabalho \cite{linuxkernel2024}.

% ----------------------------------------------------------------
\section{Segurança de contêineres}
% ----------------------------------------------------------------

Contêineres não são máquinas virtuais: compartilham o kernel do sistema hospedeiro
e dependem de mecanismos de isolamento do próprio Linux — namespaces, cgroups e
capacidades — para separar cargas de trabalho \cite{linuxkernel2024}.
Essa arquitetura torna o kernel o ponto central de controle e, consequentemente,
o principal vetor de ataque em ambientes de contêineres \cite{gregg2020}.

As ameaças mais relevantes nesse contexto são:

\textbf{Fuga de contêiner (\textit{container escape}).}
Um processo dentro de um contêiner que consegue sair do seu namespace de PID,
de rede ou de sistema de arquivos e interagir com o kernel ou com outros contêineres
sem restrição \cite{linuxkernel2024}.
Fugas exploram tipicamente vulnerabilidades em chamadas de sistema como
\texttt{unshare(2)}, \texttt{setns(2)} ou \texttt{mount(2)}, ou abusam de
capacidades Linux concedidas desnecessariamente ao contêiner.

\textbf{Escalada de privilégios dentro do contêiner.}
Mesmo sem sair do namespace, um processo que obtém \texttt{CAP\_SYS\_ADMIN}
ou executa como \texttt{root} no contêiner pode comprometer outros processos do
mesmo nó \cite{gregg2019}.

\textbf{Ataques de movimentação lateral.}
Um contêiner comprometido que estabelece comunicação direta com outros contêineres
ou serviços de infraestrutura — banco de dados, API do Kubernetes — fora dos
fluxos previstos \cite{cassagnes2020}.

O eBPF endereça esses vetores de múltiplas formas.
Programas conectados a \textit{tracepoints} de chamadas de sistema como
\texttt{unshare}, \texttt{setns} e \texttt{clone} detectam tentativas de
manipulação de namespaces em tempo real \cite{linuxkernel2024}.
\textit{Hooks} LSM BPF podem bloquear essas operações para processos cujos cgroups
não estejam em uma lista de permissões, sem qualquer modificação nas imagens de
contêiner \cite{gregg2019}.
Mapas BPF indexados por par de cgroups registram todo fluxo de rede entre
contêineres, tornando a movimentação lateral imediatamente visível \cite{majors2022}.

% ----------------------------------------------------------------
\section{Segurança em Kubernetes}
% ----------------------------------------------------------------

O Kubernetes introduz uma camada adicional de complexidade: além dos processos
em execução nos nós, existe o plano de controle — API Server, etcd, scheduler,
controller manager — cujo comprometimento implica controle total sobre o cluster
\cite{beyer2016}.
O eBPF é utilizado em duas frentes complementares nesse ambiente.

\textbf{Substituição do kube-proxy.}
O componente tradicional \texttt{kube-proxy} implementa as regras de roteamento
de serviços Kubernetes por meio de tabelas \textit{iptables} de crescimento
linear com o número de serviços \cite{linuxkernel2024}.
Projetos como o Cilium substituem o \texttt{kube-proxy} por programas eBPF anexados
a \textit{hooks} TC e XDP, implementando o mesmo roteamento com complexidade
constante e menor latência, além de oferecer visibilidade sobre cada fluxo de
rede entre \textit{pods} \cite{cassagnes2020}.

\textbf{Observabilidade do plano de dados.}
Em um cluster com centenas de \textit{pods} espalhados por dezenas de nós, rastrear
um fluxo de requisição de ponta a ponta requer correlação de eventos de múltiplos
processos e nós \cite{majors2022}.
Programas eBPF nos nós coletam identificadores de fluxo — endereços IP, portas,
metadados de cgroup — e os associam a metadados Kubernetes (namespace, nome do
\textit{pod}, \textit{labels}) obtidos do API Server, compondo uma visão de
observabilidade de rede enriquecida sem necessidade de \textit{sidecars} ou
modificação das aplicações \cite{linuxkernel2024}.

\textbf{Auditoria do API Server.}
Toda operação no cluster — criação de \textit{pods}, modificação de segredos,
concessão de permissões RBAC — passa pelo API Server e pode ser auditada por
programas eBPF conectados aos \textit{tracepoints} de rede do nó onde ele executa,
complementando os logs de auditoria nativos do Kubernetes com contexto de processo
e identidade de rede \cite{gregg2019}.

% ----------------------------------------------------------------
\section{Políticas de segurança}
% ----------------------------------------------------------------

O Kubernetes oferece primitivos nativos de política de segurança —
\textit{NetworkPolicies}, \textit{PodSecurityAdmission} e RBAC — mas sua
implementação e granularidade dependem dos plugins de infraestrutura instalados
\cite{linuxkernel2024}.
O eBPF amplia significativamente o espaço de políticas aplicáveis.

\textbf{Políticas de rede por identidade.}
Enquanto \textit{NetworkPolicies} nativas operam sobre endereços IP e portas, que
mudam a cada recriação de \textit{pod}, soluções baseadas em eBPF como o Cilium
implementam políticas baseadas em identidade: a identidade de um \textit{pod} é
derivada dos seus \textit{labels} Kubernetes e codificada em um rótulo numérico
propagado em cabeçalhos de pacotes \cite{cassagnes2020}.
Isso permite expressar regras como ``\textit{pods} do serviço A só podem receber
tráfego de \textit{pods} do serviço B no mesmo namespace'' de forma estável,
independentemente de endereços IP efêmeros \cite{linuxkernel2024}.

\textbf{Políticas de chamadas de sistema por \textit{pod}.}
Perfis seccomp restringem as chamadas de sistema disponíveis para um contêiner,
mas são estáticos e difíceis de calibrar sem causar falsos positivos \cite{gregg2020}.
O eBPF permite gerar perfis dinâmicos observando o comportamento real de cada
\textit{pod} durante um período de aprendizado e, em seguida, aplicar uma política
derivada desse perfil sem modificar o manifesto do \textit{pod}
\cite{nakryiko2020}.

\textbf{\textit{TracingPolicies} declarativas.}
Ferramentas como o Tetragon expõem um recurso Kubernetes
(\texttt{TracingPolicy}) que permite expressar políticas de segurança em YAML —
quais chamadas de sistema monitorar, quais argumentos filtrar e qual ação tomar
(alertar, registrar ou bloquear) — que são traduzidas automaticamente em programas
eBPF e implantadas em todos os nós do cluster \cite{linuxkernel2024}.
Esse modelo declarativo integra a segurança ao fluxo de \textit{GitOps} já adotado
por equipes de plataforma, tratando políticas de segurança como código versionável
e revisável \cite{beyer2016}.

% ----------------------------------------------------------------
\section{Detecção de atividades suspeitas}
% ----------------------------------------------------------------

A natureza dinâmica do Kubernetes — onde \textit{pods} são criados e destruídos
em segundos — exige que a detecção de atividades suspeitas opere em tempo real e
seja capaz de correlacionar eventos efêmeros antes que evidências desapareçam
\cite{majors2022}.
O eBPF é a tecnologia que melhor satisfaz esse requisito, e os padrões de ataque
mais comuns em ambientes Kubernetes têm contrapartidas diretas em programas de
detecção baseados em eBPF.

\textbf{Execução de shell em contêiner.}
A abertura de um terminal interativo (\texttt{/bin/sh}, \texttt{/bin/bash})
dentro de um contêiner de produção é um forte indicador de comprometimento ou
de acesso não autorizado via \texttt{kubectl exec} \cite{cassagnes2020}.
Programas eBPF detectam esse padrão combinando o \textit{tracepoint}
\texttt{sys\_enter\_execve} com a verificação de que o processo possui um TTY
associado e que seu cgroup corresponde a um \textit{pod} em execução
\cite{linuxkernel2024}.

\textbf{Acesso a segredos montados.}
Segredos Kubernetes são montados como arquivos no sistema de arquivos do contêiner.
Leitura desses arquivos por processos diferentes do esperado — especialmente por
shells ou binários de sistema — é detectada por programas eBPF conectados ao
\textit{tracepoint} \texttt{sys\_enter\_openat}, filtrando caminhos que
correspondam a diretórios de montagem de segredos (\texttt{/var/run/secrets/})
\cite{gregg2019}.

\textbf{Comunicação com o API Server a partir de contêineres de aplicação.}
Requisições HTTP ao endereço do API Server (\texttt{kubernetes.default.svc})
originadas de \textit{pods} de aplicação que não deveriam ter acesso ao plano
de controle indicam possível tentativa de movimentação para comprometer o cluster
\cite{beyer2016}.
Mapas BPF indexados por endereço de destino e cgroup de origem permitem detectar
esses fluxos em microssegundos e bloqueá-los via \textit{hook} de cgroup de socket
\cite{linuxkernel2024}.

\textbf{Modificação de imagens em tempo de execução.}
Contêineres devem ser imutáveis: qualquer escrita no sistema de arquivos raiz
fora de volumes montados explicitamente é suspeita \cite{gregg2020}.
Programas eBPF conectados a \textit{hooks} LSM de abertura de arquivo em modo de
escrita — filtrados pelo namespace de montagem do contêiner — detectam esse padrão
independentemente da ferramenta ou técnica usada pelo atacante para realizar a
escrita \cite{linuxkernel2024}.

\textbf{Varredura de rede interna.}
Um \textit{pod} comprometido que tenta conectar a múltiplos endereços e portas
em sequência rápida executa varredura de rede para mapear serviços internos do
cluster \cite{cassagnes2020}.
Mapas BPF do tipo \texttt{BPF\_MAP\_TYPE\_LRU\_HASH} acumulam contagens de
tentativas de conexão por \textit{pod} em janelas de tempo deslizantes e acionam
alertas quando o número de destinos distintos excede um limiar configurável
\cite{linuxkernel2024,gregg2020}.